% Pursuing Stacks (cotes 134-2 to 134-8): archive cote 134-2, site page 31 (PDF page 32).
% letter to Quillen. Draft compared: ps1.tex:823-894.
% AI-transcribed by Claude Opus 5.5 (claude-opus-5-5, adaptive thinking, effort high), Message Batches API, batch ps1 (msgbatch_01Rv5hCaNTtUt4JHx1yqoHPL), 2026-11-12; not yet checked by eye.
\scanpage{134-2/31}
\mspage{9'}
infinitely many ope\xout{x}rations standing for the single ope\xout{\unsure{rat}}ion $v\overset{i}{\circ}u$ say,
which could be \unsure{wiewed|viewed} in effect in terms of an arbitrary n-sequence
($n\geqslant 2$) of "composable" i-objects $U_1=u,u_2=v,\ u_3.\dots\ u_n$ ;\inserted{just ``forgetting'' $u_3,\dots,u_n$!} The natural
way to meet this "objection" would be to restrict to \xout{a} pairs (f,g)
which cannot be factored non trivially through another object $K'$ as
\diagram{\begin{tikzcd} K \ar[r] & K' \ar[r, shift left, "{f',g'}"] \ar[r, shift right] & F_i \end{tikzcd} .}
But even with such restrictions, there remain a lot of redundancies
- and this again seems to me in the nature of things, namely that there
is no really natural, "most economic" way for achieving condition (A),
by a stepwise construction meeting condition (B). For instance, in $\underline{C}_1$
already we will have \xout{the} not merely the compositions $v\overset{i}{\circ}u$ , but at the
same time simultaneous compositions
\[ (*) \qquad u_n\circ u_{n-1}\circ - - \circ u_1 \]
for "composable" sequences of i-objects ($i\geqslant 1$), without reducing this
(as is customary) to iteration of the binary composition $v\circ u$ . Of course
using the binary composition, and more generally iteration of $n'$-ary
compositions with $n'<n$ (when $n\geqslant 3$), we get an impressive bunch of
ope\xout{x}rations in the n variables $u_1$ , ..., $u_n$ , serving the same purpose
as (*). All these will be tied up by homotopies in the next step $\underline{C}_2$ .
We would like to think of this set of homotopies in $\underline{C}_2$ as a kind of
"transitive system of isomorphisms" (of associativity); now the transiti%
% NOTE: the right edge of the scan cuts the end of this line and of some following lines
vity relations one is looking for will be replaced by homotopies again
between compositions of homotopies, which will enter in the picture with
$\underline{C}_3$ ,\inserted{etc.} Here the infinite story is exemplified by the more familiar situa-%
\margin{$\otimes$} tion of the two ways in which one could define a\hand{"}\hand{$\otimes$}-composition with
associativity" in a category, \xout{\illegible{1}} starting \uline{either} in terms of
a binary operation, \uline{or}\inserted{with} a bunch of n-ary operations - with\hand{,} in both \deleted{cas\unsure{e}}\inserted{cases,}
the associativity isomorphisms being an essential part of the structure.
Here again, while it is generally (and quite validly) felt that the
two points of view are\hand{"}equivalent\hand{"}, and both have their advantages and
their drawbacks, still it is not true\inserted{I Believe} that the \inserted{two} categories of algebraic
\margin{$\otimes$} structures "category with associative \hand{$\otimes$}-operation", using one or the
\margin{* even not in the \unsure{grou-poid} context} other definition, are equivalent\hand{*}.\inserted{(15)} Here the story though of the relation
between the two notions is a \uline{finite} one, due to the fact that it is
related to the notion of 2-categories or 2-groupoids, instead of $\infty$ -
groupoids as before ...

Thus I do\hand{'}nt feel really like spending much energy in cutting down
redundancies, but prefer working with a notion of $\infty$-groupoid which \deleted{remain\unsure{s}}\inserted{remains}
partly indeterminate, the main features being embodied in the condition\unsure{\hand{s}}
(A) and (B) and in the description of $\underline{C}_0$ , without other specification.

\margin{31}
